\documentclass{article}
\usepackage[UTF8]{ctex}
\usepackage[a4paper,
    left=1.75cm,
    right=1.75cm,
    top=2.25cm,
    bottom=2.25cm
]{geometry}
\usepackage{amsmath, amssymb, amsthm}
\usepackage{xcolor}
\usepackage{graphicx}
\usepackage{enumerate}
\usepackage[normalem]{ulem}
\usepackage{subcaption}
\usepackage{fancyhdr}

\newcommand{\highlight}[1]{\textcolor{blue}{#1}}

\pagestyle{fancy}
\fancyhf{} % 清空默认
\fancyfoot[R]{\thepage}

% Problem 环境
\newtheorem{definition}{Definition}
\newtheorem{problem}{Problem}           % 如果要中文可以改成 {问题}
\newtheorem{proposition}{Proposition}
\newtheorem{example}{Example}
\newtheorem{assumption}{Assumption}

\def\highlight#1{\textcolor{blue}{#1}}
\def\hhighlight#1{\textcolor[rgb]{0.765, 0.682, 0.796}{#1}}
\def\corrected#1{\textcolor{cyan}{#1}}

\title{Brief Discussion on cognitive non-blocking}
\author{-}
\date{April 2026}

\begin{document}

\maketitle

\begin{enumerate}[(i)]
    \item What is cognitive non-blocking in DES? How to formulate the problem?
    \item What is the assumption for the system?
    \item (Optional) How to solve cognitive non-blocking?
\end{enumerate}

% Preliminaries
\highlight{
We say a FSA is non-block if and only if
$\forall s \in \mathcal{L}(S/G), \exists t \in \Sigma^*, \delta(st) \subseteq X_m$}

Given a closed-loop system $S/G$, We define the sublanguage ending with an observation as
$\mathcal{L}_{op}(S/G) = \overline{(\Sigma^* \Sigma_o) \cap \mathcal{L}(S/G)}$.
We say that the supervisor $S$ is \textit{cognitive non-blocking} if
$\mathcal{L}(S / G) = \overline{\mathcal{L}_{op}(S/G) \cap \mathcal{L}_m(G)}$.
Intuitively, a closed-loop system is cognitive non-blocking if, from every reachable observation, the system can generate some further observations that certify that the marked set $X_m$ has been visited.

We define
$$
\hat{\mathcal{E}}_S(s) = \{ x \in X | \exists t \in \Sigma^*\Sigma_{uo}^*, \sigma_{uo} \in \Sigma_o \ s.t. \ t \sigma_{uo} \in \mathcal{L}(S / G) \wedge \delta(x_0, t) = x \wedge P(t \sigma_{uo}) =  s \}
$$
as the \textit{cognitive state estimate} upon the occurrence of $s$, which is the set of states the system could be in before the final observation of $s$.

Then a closed-loop system $S / G$ is cognitive non-blocking if and only if
$$
\forall r \in \mathcal{L}(S / G), \exists t (\in \Sigma^*\Sigma_o), \hat{\mathcal{E}}_S(P(rt)) \subseteq X_m
$$

\begin{definition}[Cognitive Non-Blocking]
    TODO
\end{definition}


\end{document}